\section{Problem Setting and Benchmark Definition}
\label{sec:problem}

We evaluate early-warning detectors on simulated trajectory data from three
dynamical systems. Each detector receives a trajectory and produces a score
sequence. The score at time \(t\), denoted \(p_{i,t}\), is interpreted as the
detector's warning score for trajectory \(i\). Higher values indicate stronger
evidence for an approaching transition.

The benchmark is defined to compare observer behavior across regimes rather
than to select a single global winner. All reported comparisons use the same
systems, patient-count settings, seed order, thresholding rule, and evaluation
metrics.

\subsection{Dynamical Systems}
\label{sec:systems}

The benchmark contains three synthetic systems.

\begin{description}[leftmargin=*]
  \item[Fold.]
  The fold system represents a collapse-like transition in which a stable state
  is lost as the control condition changes. It tests whether an observer can
  detect gradual loss of stability without relying on oscillatory structure.

  \item[Hopf.]
  The Hopf system represents an oscillatory instability. It tests whether an
  observer can use temporal structure associated with changing stability and
  emerging oscillatory behavior.

  \item[Logistic.]
  The logistic system represents a discrete-time transition with
  period-doubling behavior. It tests whether an observer calibrated on
  continuous warning scores also behaves well for map-like dynamics.
\end{description}

These systems are not treated as interchangeable datasets. They are distinct
test cases for the same observer family.

\subsection{Trajectory Data and Patient Counts}
\label{sec:data}

For each system, the benchmark contains signal trajectories and null
trajectories. Signal trajectories contain a transition with a bifurcation time
\(\tau_i\). Null trajectories provide non-transition examples for
false-positive evaluation.

The patient count denotes the number of simulated trajectories used in a
benchmark condition. The finalized patient-count settings are
\[
  n \in \{100, 200, 300, 400, 500\}.
\]
For each patient count, system, and observer variant, results are aggregated
over the following ten seeds:
\[
  101, 202, 303, 404, 505, 606, 707, 808, 909, 1010.
\]

The finalized reporting pipeline compares two learned observer variants:
\(\text{Kalman-BCE}\) and \(\text{Kalman-LSTM-Spec}\). Each method produces
test scores \(p_{i,t} \in [0,1]\), which are used for threshold selection,
detection-time computation, early-warning AUC, and false-positive evaluation.

\subsection{Batch Selection}
\label{sec:batch}

The experiment directory can contain multiple batches for the same patient
count, including partial reruns. To avoid manual selection, the analysis uses a
deterministic batch-selection rule. For each patient count \(n\), let
\(\mathcal{B}_n\) be the set of available batches. The selected batch is the
batch that maximizes, in order:

\begin{enumerate}[leftmargin=*]
  \item the number of unique \((\text{system}, \text{method}, \text{seed})\)
        triples,
  \item the number of systems represented,
  \item the total number of result records,
  \item and the timestamp, if the preceding quantities are tied.
\end{enumerate}

This rule keeps the most complete batch for each patient count and excludes
partial reruns from the main figures and tables.

\subsection{Evaluation Metrics}
\label{sec:metrics}

The benchmark reports three metrics: detection time, early-warning AUC, and
false-positive rate. The metrics measure different properties of the same score
sequence and should not be collapsed into a single ordering.

\subsubsection{Detection Time}
\label{sec:metrics:dt}

For a positive trajectory \(i\) with bifurcation time \(\tau_i\), define the
first pre-bifurcation alert time as
\[
  a_i = \min \{t : 0 \leq t < \tau_i,\; p_{i,t} \geq \theta \},
\]
where \(\theta\) is the decision threshold. If such an alert exists, the
trajectory-level detection time is
\[
  DT_i = \tau_i - a_i.
\]
The reported detection time is the mean of \(DT_i\) over positive trajectories
with at least one pre-bifurcation alert. Trajectories without a pre-bifurcation
alert are omitted from this mean rather than counted as zero-lead detections.

\subsubsection{Early-Warning AUC}
\label{sec:metrics:auc}

For each positive trajectory, the early-warning score is the maximum model
score in a window before the bifurcation:
\[
  e_i = \max \{p_{i,t} : \max(0,\tau_i - 50) \leq t < \tau_i - 5\}.
\]
For each null trajectory \(j\) with sequence length \(T_j\), the corresponding
score is computed over the terminal part of the sequence:
\[
  n_j = \max \{p_{j,t} : \max(0,T_j - 50) \leq t < T_j - 5\}.
\]
The early-warning AUC is the ROC AUC computed from the pooled signal scores
\(\{e_i\}\) and null scores \(\{n_j\}\), with labels one for signal and zero
for null.

\subsubsection{False Positive Rate}
\label{sec:metrics:fpr}

False-positive rate is computed on null trajectories as a per-time-step alert
rate:
\[
  FPR =
  \frac{\sum_j \sum_{t=0}^{T_j-1} \mathbf{1}[p_{j,t} \geq \theta]}
       {\sum_j T_j}.
\]
This quantity measures the fraction of null time steps that would trigger an
alert. It is not a per-trajectory false-positive probability.

\subsubsection{Threshold Selection}
\label{sec:metrics:threshold}

Thresholds are selected on the validation split before test evaluation. For a
positive validation trajectory, scores in the window
\[
  \max(0,\tau_i - 60) \leq t < \tau_i
\]
are labeled positive. Scores from the early pre-bifurcation region
\[
  0 \leq t < \max(\max(0,\tau_i - 120) - 1, 0)
\]
are labeled negative. If null validation scores are supplied, their first 120
time steps are also labeled negative. In the finalized benchmark runs reported
here, threshold selection for the two learned observer variants uses the signal
validation trajectories.

The threshold is selected from the validation ROC curve by maximizing Youden's
statistic \citep{youden1950index}, \(J = TPR - FPR\). If the sensitivity at that threshold is below
0.80, the selected threshold is the one whose sensitivity is closest to 0.80.
If the validation labels contain only one class, or if the validation scores are
effectively constant, the fallback threshold is 0.5.
